\section{Experimental Setup}
\label{sec:setup}

We organize the evaluation around six questions. RQ1 asks whether \qss{} returns complete evidence sets with small budgets and what its set formulation adds over the same-model controls. RQ2 asks whether better set recall improves answers on the development benchmark at comparable cost, and how robust these gains are. RQ3 tests these conclusions on a held-out benchmark with pre-registered hypotheses. RQ4 asks where the selection capability resides and whether it transfers to lightweight selectors, to write time, or to other model families. RQ5 asks whether the failures diagnosed on held-out data can be turned into a mechanism that improves answers on further held-out data. RQ6 asks whether set selection can approach the cost of retrieval with a lightweight reader and beat retrieval with a strong reader at no higher cost. Appendix~\ref{app:zoom} reports two further ideas for multimodal memory, reader-driven zooming and captions as memory content, neither of which helped.

\paragraph{Development benchmark.} MemEye~\citep{guo2026memeye} provides 371 open-ended questions over eight scenario histories of 30--299 rounds and 30--89 images each, with annotated evidence rounds for every question. We designed the index, the prompts, and all lightweight variants on these questions. Mem-Gallery~\citep{bei2026memgallery} serves as a second benchmark (a stratified subset of 903 questions over its 20 scenarios). For a scale test, each MemEye history is interleaved with Mem-Gallery sessions that contribute 1,000 additional images; these sessions come from other users, are re-dated within the span of the target history, and leave the questions and their evidence unchanged, while the mean history grows from 126.5 to 4,088.5 rounds.

\paragraph{Held-out benchmark.} MemLens~\citep{ren2026memlens} contains multi-session conversations with images in which every question has its own haystack; its authors report that removing the evidence images drops frontier models below 2\% accuracy on the questions whose evidence includes images. We use all 699 questions except the abstention type at the 32k context length (64.6 rounds and 13.3 images per question) and a fixed random subset of 200 questions at 128k (245 rounds and 48.8 images per question). A round is a user turn with the following assistant turn, and the gold rounds are the rounds of the annotated answer sessions that contain an evidence turn or an evidence image, 5.9 per question on average. Before running any condition, we froze the index, prompts, and parameters from MemEye, registered the conditions and five hypotheses, and declared one variant: index lines that keep 600 instead of 200 characters of the user message, motivated only by the length of the evidence turns (median 637 characters). MemLens played no role in any earlier design decision. For the mechanism study of RQ5, which the outcome of this test motivated, we split the 699 questions into a development third (233 questions) and a test part (466 questions), stratified by question type and subtype; we developed all new mechanisms on MemEye and the MemLens development part only, and we used the test part once, for conditions and hypotheses registered beforehand. The deployment mechanisms of RQ6 were developed on the same data; their registered confirmation uses the test-split questions with the 64k haystacks of MemLens (124.8 rounds and 25.5 images per question), which are twice as long as those seen before, and the persistent histories of Mem-Gallery.

\paragraph{External benchmark.} To test the frozen VGS design outside MemLens, we use the lifelong-personal part of WorldMemArena~\citep{liu2026worldmemarena}, which neither our design nor any earlier decision touched: 20 user histories of 30 sessions with shared images, whose 1,100 questions are asked at 100 checkpoints after growing prefixes of the sessions (146.7 rounds per history on average, median user message of 397 characters). The benchmark annotates evidence as session-level memory points rather than rounds, so we evaluate this test end to end. We registered the conditions and hypotheses before any run; the registration originally included the lifelong-project part, which we dropped before any run because its images could not be downloaded in time. Everything else is frozen, including our captions, the index, the prompts, the visibility estimate, and the threshold $\tau{=}0.7$.

\paragraph{Models.} The lightweight reader is \texttt{gemini-2.5-flash-lite} at temperature 0 with images at its default resolution, run without thinking and, where stated, with a 1,024-token thinking budget; without thinking it decodes deterministically. The strong reader is \texttt{gemini-3.8-flash} with low thinking and low-resolution images. The \qss{} selector is \texttt{gemini-3.8-flash} with low thinking; lightweight selectors, the router of set-aware routing, and the note writer of VGS use \texttt{gemini-2.5-flash-lite} without thinking. The strong selector does not decode deterministically, so we report one selection draw per condition unless stated otherwise and quantify the variation across four draws in RQ6. To test transfer across model families, we run the same selector prompt with \texttt{gpt-5}, \texttt{gpt-5-mini}, and Qwen3-235B-A22B in its instruction and thinking versions. Image descriptions are written once by \texttt{gemini-3.8-flash} from the image alone, and embeddings use \texttt{gemini-embedding-2} with a key of round text plus caption.

\paragraph{Baselines.} Full-context reading (FC) gives the reader the entire history with all images. Top-$k$ retrieval (RAG@$k$) gives the reader the $k$ rounds with the highest cosine similarity, in chronological order and with the original images. Annotated evidence gives the reader exactly the gold rounds; it is a reference for focused reading rather than an upper bound. All conditions with the same reader share one prompt, and \qss{} differs from RAG@$k$ and from the two controls only in how the rounds are chosen.

\paragraph{Judges and statistics.} On MemEye, the official 0--1 rubric is applied by \texttt{gemini-3.8-flash}, and an independent judge from another family (Qwen3-VL-235B-A22B) re-scores the main conditions. On MemLens, we use the official judge prompts with Qwen3-VL-235B-A22B, the judge family of the benchmark authors, and re-score the main conditions with \texttt{gemini-3.8-flash}. Retrieval is measured by full-set recall, coverage, and the number of opened rounds. Cost is the measured spend per 1,000 questions at list prices over all calls except the judge, including the creation and storage of context caches where they are used. Differences are paired over the same questions with 10,000 bootstrap resamples. Because MemEye questions share eight histories, we also report a two-stage bootstrap that resamples scenarios and then questions within them; MemLens questions have independent haystacks.
